% ============================================================
%  ch09 —— §5.6（原文 p.37–p.40）
% ============================================================
\routenote{p.37--p.40}

\subsection{定理 4(2) 的证明}

本节假设 $G$ 型为 $B,C$ 或 $F_4$。对型 $A,D,E$ 的 $G$，定理 4 的陈述 (2) 与 (1) 相同；$G=G_2$ 的情形将在 §11 中证明。

把推论 5.20 用于 $V=\ckg$（它有两个非零支配权 $\theta^\vee$ 与 $\gamma^\vee$，重数均为 $1$），得
\[
\chi_c(\PP^1\setminus\{0,\infty\},\Kl^{\mathrm{Ad}}_{\ckG}(\varphi,\chi))=\chi_c\big(\mathrm{Gr}^{\mathrm{triv},a=0}_{\theta^\vee},\mathrm{ev}^*J\big)+\chi_c\big(\mathrm{Gr}^{\mathrm{triv},a=0}_{\gamma^\vee},\mathrm{ev}^*J\big).
\]
由 (1) 已知 $\chi_c(\mathrm{Gr}^{\mathrm{triv},a=0}_{\theta^\vee},\mathrm{ev}^*J)=r_\ell(G)$，只需证明
\[
\chi_c\big(\mathrm{Gr}^{\mathrm{triv},a=0}_{\gamma^\vee},\mathrm{ev}^*J\big)=r_s(G).
\]

对任意子集 $S\subset\Phi$，令
\[
V^S_\beta=\Big\{\prod u_{\beta''}(c_{\beta''}\tau^{-1})\in V_\beta\ \Big|\ c_{\beta''}=0\ \text{对}\ \beta''\notin S\Big\}.
\]
利用此定义与 $\Phi$ 上的标准偏序（由注 5.34(1)，它限制为 $\Phi^\beta_2$ 上的全序），$V^{>-\theta}_\beta$、$V^{>0}_\beta$ 等记号的含义是显然的。

利用引理 5.31、引理 5.35 与分解 (5.32)，消解 $\nu$ 限制到 $\mathrm{Gr}^{\mathrm{triv},a=0}_{\le\gamma^\vee}$ 上可读作
\[
\nu_{a=0}:\bigsqcup_{\beta\ \text{短根}}\mathrm{Ad}(U)\big(V^{>-\theta}_\beta\big)\to\mathrm{Gr}^{\mathrm{triv},a=0}_{\le\gamma^\vee}.
\]
由引理 5.31，$\mathrm{ev}^*J$ 在 $\nu_{a=0}$ 的目标上的 Euler 特征有来自 $\star$、$\mathrm{Gr}^{\mathrm{triv},a=0}_{\theta^\vee}$ 与 $\mathrm{Gr}^{\mathrm{triv},a=0}_{\gamma^\vee}$ 的贡献：
\begin{equation}
\begin{aligned}
\sum_{\beta\ \text{短根}}\chi_c\big(\mathrm{Ad}(U)V^{>-\theta}_\beta,\mathrm{ev}^*J\big)
&=\chi_c(G/P_{-\gamma})+\chi_c\big(\mathrm{Gr}^{\mathrm{triv},a=0}_{\theta^\vee},\mathrm{ev}^*J\big)\chi_c(L_\theta/L_\theta\cap P_{-\gamma})\\
&\quad+\chi_c\big(\mathrm{Gr}^{\mathrm{triv},a=0}_{\gamma^\vee},\mathrm{ev}^*J\big).
\end{aligned}
\end{equation}
Bruhat 分解给出
\begin{equation}
\chi_c(G/P_{-\gamma})=\#(W/W_\gamma)=\#\{\Phi\ \text{中的短根}\};
\end{equation}
\begin{equation}
\chi_c(L_\theta/L_\theta\cap P_{-\gamma})=\#(W_\theta/W_\theta\cap W_\gamma).
\end{equation}

\medskip\noindent\textbf{断言.} 对短根 $\beta$，
\[
\chi_c\big(\mathrm{Ad}(U)(V^{>-\theta}_\beta),\mathrm{ev}^*J\big)=
\begin{cases}
0,&\Phi^\beta_2\ \text{含单根},\\
1,&\text{否则}.
\end{cases}
\]

先承认该断言来完成证明。结合 (5.36)、(5.37)、(5.38) 与断言，得
\begin{equation}
\begin{aligned}
-\chi_c\big(\mathrm{Gr}^{\mathrm{triv},a=0}_{\gamma^\vee},\mathrm{ev}^*J\big)
&=\#\{\text{短根}\}-r_\ell(G)\#(W_\theta/W_\theta\cap W_\gamma)\\
&\quad-\#\{\text{短根}\ \beta\mid\Phi^\beta_2\ \text{不含单根}\}\\
&=\#\{\text{短根}\ \beta\mid\Phi^\beta_2\ \text{含单根}\}-r_\ell(G)\#(W_\theta/W_\theta\cap W_\gamma).
\end{aligned}
\end{equation}
由注 5.34(1)，$\Phi^\beta_2$ 按高度全序，故至多含一个单根。因此
\[
\#\{\text{短根}\ \beta\mid\Phi^\beta_2\ \text{含单根}\}=N_s+N_\ell,
\]
其中 $N_s$（相应地 $N_\ell$）是使 $\Phi^\beta_2$ 含短（相应地长）单根的短根 $\beta$ 的个数。
\begin{itemize}[leftmargin=1.6em, itemsep=.2em]
\item 若 $\Phi^\beta_2$ 含短单根，则由注 5.34(2)，该单根必为 $\beta$，因为 $\beta$ 是 $\Phi^\beta_2$ 中唯一的短根。故 $N_s=r_s(G)$。
\item 任一长单根都在 $\theta$ 的 $W$-轨道中，故 $N_\ell=r_\ell\cdot N_\theta$，其中 $N_\theta=\#\{\text{短根}\ \beta\mid\theta\in\Phi^\beta_2\}$。这样的短根 $\beta$ 在 $W_\theta$-轨道中（这由把注 5.34(2) 用于对偶根系得到），而 $\gamma$ 是其中之一，故 $N_\theta=\#(W_\theta/W_\theta\cap W_\gamma)$。于是 $N_\ell=r_\ell(G)\#(W_\theta/W_\theta\cap W_\gamma)$。
\end{itemize}
结合这些计算与 (5.39)，即得 $-\chi_c(\mathrm{Gr}^{\mathrm{triv},a=0}_{\gamma^\vee},\mathrm{ev}^*J)=r_s(G)$，从而证明定理 4(2)。

剩下证明断言。

\emph{情形 I：$\Phi^\beta_2$ 含单根 $\alpha_i$。} 对 $u\in U$ 与 $v=u_{\alpha_i}(c_{\alpha_i}\tau^{-1})v_{\alpha_i}\in V_\beta$（其中 $v_{\alpha_i}\in\prod_{\beta'\neq\alpha_i}U_{\beta',1}$），有
\[
\mathrm{ev}(\mathrm{Ad}(u)v)=\mathrm{Ad}(u)u_{\alpha_i}(c_{\alpha_i})\cdot\mathrm{Ad}(u)\mathrm{ev}(v_{\alpha_i}).
\]
这意味着 $\mathrm{Ad}(u)v\in UB^{\mathrm{opp}}$ 当且仅当 $\mathrm{Ad}(u)\mathrm{ev}(v_{\alpha_i})\in UB^{\mathrm{opp}}$；且当这发生时，它在 $U/[U,U]$ 中的像是 $u_{\alpha_i}(c_{\alpha_i})$ 乘以 $\mathrm{Ad}(u)\mathrm{ev}(v_{\alpha_i})$ 在 $U/[U,U]$ 中的像。

利用分解 (5.33)，可把 $\mathrm{Ad}(U)V^{>-\theta}_\beta$ 写成 $U_{\alpha_i,-1}\times A$，其中 $A$ 是某个仿射空间（其余因子之积），且 $\mathrm{ev}^*J$ 是外积 $\mathrm{ev}^*\mathrm{AS}_{\varphi_i}\boxtimes L$，$L$ 是 $A$ 上某局部系统。由 Künneth 公式，
\[
H^*_c\big(\mathrm{Ad}(U)V^{>-\theta}_\beta,\mathrm{ev}^*J\big)=H^*_c\big(U_{-\alpha_i,-1},\mathrm{ev}^*\mathrm{AS}_{\varphi_i}\big)\otimes H^*_c(A,L)=0.
\]

\emph{情形 II：$\Phi^\beta_2$ 不含任何单根。} 把向量空间 $V^{>-\theta}_\beta$ 分层为
\[
V^{>-\theta}_\beta=V^{>0}_\beta\ \sqcup\ \Big(\bigsqcup_{\beta'\in\Phi^\beta_2,\ -\theta<\beta'<0}\big(V^{\ge\beta'}_\beta-V^{>\beta'}_\beta\big)\Big).
\]
首先证明 $\chi_c(\mathrm{Ad}(U)V^{>0}_\beta,\mathrm{ev}^*J)=1$。事实上，由于 $\Phi^\beta_2\cap\Phi^+$ 不含单根，$\mathrm{ev}(\mathrm{Ad}(U)V^{>0}_\beta)\subset[U,U]$，故 $\mathrm{ev}^*J$ 是 $\mathrm{Ad}(U)V^{>0}_\beta$ 上的常层。又由 (5.33)，$\mathrm{Ad}(U)V^{>0}_\beta$ 是仿射空间，故得结论。

其次证明对每个 $\beta'\in\Phi^\beta_2$、$-\theta<\beta'<0$，有 $\chi_c(\mathrm{Ad}(U)(V^{\ge\beta'}_\beta-V^{>\beta'}_\beta),\mathrm{ev}^*J)=0$。由于 $\beta'\neq-\theta$，存在单根 $\alpha_i$ 使 $\beta'-\alpha_i$ 仍是根。则 $\alpha=\alpha_i-\beta'$ 是正根。按假设 $\alpha_i\notin\Phi^\beta_2$，故 $\langle\alpha,\beta^\vee\rangle=\langle\alpha_i,\beta^\vee\rangle-2<0$，即 $\alpha$ 出现在 (5.33) 的第一个乘积中。

\begin{lemma}
设 $a,b\in\mathbb{Z}_{>0}$，$\beta'\le\beta_1,\dots,\beta_b$ 是 $\Phi^\beta_2$ 中的根（不必互异），则 $a\alpha+\sum_{i=1}^b\beta_i\neq 0$。若 $a\alpha+\sum_{i=1}^b\beta_i$ 是根，则下列情形之一发生：
\begin{enumerate}[label=(\arabic*), leftmargin=1.9em, itemsep=.2em]
\item $a\alpha+\sum_{i=1}^b\beta_i$ 是负根。此时 $a=b-1$，$a\alpha+\sum_{i=1}^b\beta_i\in\Phi^\beta_2$，且（在 $\Phi^\beta_2$ 的全序中）大于任一 $\beta_i$；
\item $a\alpha+\sum_{i=1}^b\beta_i$ 为正但非单根；
\item $a=b=1$，$\beta_1=\beta'$，且 $a\alpha+\sum_{i=1}^b\beta_i=\alpha_i$。
\end{enumerate}
\end{lemma}

\begin{proof}
若 $a\ge b$。由于 $\alpha+\beta'>0$ 且 $\beta_i\ge\beta'$，$\alpha+\beta_i$ 必有正高度。因此
\[
a\alpha+\sum_{i=1}^b\beta_i=(a-b)\alpha+\sum_{i=1}^b(\alpha+\beta_i)>0.
\]
由于上式右端每一项都有正高度，当只有一个被加项时它必是单根，而这只能发生于情形 (3)。

若 $a<b$，则 $\langle a\alpha+\sum_{i=1}^b\beta_i,\beta^\vee\rangle\ge 2b-2a\ge 2$。因此若 $a\alpha+\sum_{i=1}^b\beta_i$ 是根，必有 $a=b-1$ 且 $a\alpha+\sum_{i=1}^b\beta_i\in\Phi^\beta_2$。由于 $\Phi^\beta_2$ 不含单根，得情形 (1) 或 (2)。

无论如何有 $a\alpha+\sum_{i=1}^b\beta_i\neq 0$。
\end{proof}

对 $u=\prod u_{\alpha'}(c_{\alpha'})$（乘积跑遍 $\alpha'\in\Phi^+$、$\langle\alpha,\beta^\vee\rangle<0$），记 $u_{\underline\alpha}=\prod_{\alpha'\neq\alpha}u_{\alpha'}(c_{\alpha'})$。类似地，对 $v\in V^{\ge\beta'}_\beta-V^{>\beta'}_\beta$，写 $v=u_{\beta'}(c_{\beta'}\tau^{-1})v_{\beta'}\prod u_{\beta''}(c_{\beta''}\tau^{-1})$（乘积跑遍 $\beta'<\beta''\in\Phi^\beta_2$，$c_{\beta'}$ 可逆）。则
\begin{equation}
\mathrm{ev}(\mathrm{Ad}(u)v)=u_\alpha\big[u_{\underline\alpha}(c_\alpha),u_{\beta'}(c_{\beta'})\big]u_{\beta'}(c_{\beta'})\Big(\prod_{\beta''>\beta'}\big[u_{\underline\alpha}(c_\alpha),u_{\beta''}(c_{\beta''})\big]u_{\beta''}(c_{\beta''})\Big)u_{\underline\alpha}^{-1}.
\end{equation}
为计算 (5.41)，用 Chevalley 交换子关系把每个 $[u_{\underline\alpha}(c_\alpha),u_{\beta''}(c_{\beta''})]$ 写成正根与负根因子之积，并设法把负根因子（由引理 5.40(1)，它们必出现在 $\Phi^\beta_2$ 中）移到右边。由 Chevalley 交换子关系，每一步都会产生形如 $a\alpha+\beta_1+\cdots+\beta_b$（$\beta_i\in\Phi^\beta_2$，$a,b>0$）的新因子。我们保留正因子，把负因子继续往右移。此过程中唯一需保证的是：在做交换子 $[U_{\alpha'},U_{\gamma'}]$ 时总有 $\alpha'$ 与 $\gamma'$ 线性无关，从而 Chevalley 交换子关系可用。事实上，过程中只出现 $\alpha'\in\Phi^+$、$\gamma'\in\Phi^-$ 且 $\alpha'+\gamma'$ 形如 $a\alpha+\beta_1+\cdots+\beta_b$（$\beta'\le\beta_i\in\Phi^\beta_2$，$a,b\in\mathbb{Z}_{>0}$）的情形。$\alpha'$ 与 $\gamma'$ 线性相关的唯一可能是 $\alpha'+\gamma'=0$，而这已被引理 5.40 排除。

过程结束时得
\begin{equation}
\mathrm{ev}(\mathrm{Ad}(u)v)=u_\alpha u_{\alpha_i}(\epsilon c_\alpha c_{\beta'})u_+\Big(\prod_{\beta''\ge\beta',\ \beta''\in\Phi^\beta_2}u_{\beta''}(c_{\beta''}+e c_{\beta''})\Big)u_{\underline\alpha}^{-1}.
\end{equation}
其中项 $u_{\alpha_i}(\epsilon c_\alpha c_{\beta'})$（$\epsilon=\pm 1$）来自 $[u_\alpha(c_\alpha),u_{\beta'}(c_{\beta'})]$。项 $u_+$ 是 $U_{a\alpha+\beta_1+\cdots+\beta_b}$（即 $a\alpha+\beta_1+\cdots+\beta_b\in\Phi^+$）中所有其他正因子之积。由引理 5.40(2)(3)，这些 $a\alpha+\beta_1+\cdots+\beta_b$ 从不是单根，故 $u_+\in[U,U]$。最后，额外系数 $e c_{\beta''}$ 来自 $U_{a\alpha+\beta_1+\cdots+\beta_b}$ 中的负因子，它是 $c_\alpha$ 与 $c_{\beta'''}$（$\beta'''\in\Phi^\beta_2$）的多项式函数。由引理 5.40(1)，$c_{\beta''}$ 只涉及那些满足 $\beta'''<\beta''$ 的 $c_{\beta'''}$。

因此可作变量替换
\begin{equation}
\begin{aligned}
\mathrm{Ad}(U)\big(V^{\ge\beta'}_\beta-V^{>\beta'}_\beta\big)&\xrightarrow{\ \sim\ }U_\alpha\times U^\times_{\beta',-1}\times\prod_{\langle\alpha',\beta^\vee\rangle<0,\ \alpha'\neq\alpha}U_{\alpha'}\times\prod_{\beta''\in\Phi^\beta_2,\ \beta''>\beta'}U_{\beta'',-1},\\
\mathrm{Ad}(u)(v)&\leftrightarrow\Big(u_\alpha(c_\alpha),\ u_{\beta'}(c_{\beta'}\tau^{-1}),\ u_{\underline\alpha},\ \prod u_{\beta''}(c_{\beta''}+ec_{\beta''})\Big).
\end{aligned}
\end{equation}
设 $A$ 为 (5.43) 中后三项之积，$\mathrm{pr}_A:\mathrm{Ad}(U)(V^{\ge\beta'}_\beta-V^{>\beta'}_\beta)\to A$ 为投影。由 (5.42)，$\mathrm{ev}^*J$ 限制在 $\mathrm{pr}_A$ 的纤维上同构于 $U_\alpha$ 上的 Artin--Schreier 层，故 $\mathrm{pr}_{A,!}\mathrm{ev}^*J=0$，从而
\[
H^*_c\big(\mathrm{Ad}(U)U^\times_{\beta,-1},\mathrm{ev}^*J\big)=H^*_c\big(A,\mathrm{pr}_{A,!}\mathrm{ev}^*J\big)=0.
\]
这完成了断言的情形 II。
